\documentclass[11pt]{article}
\usepackage[margin=1in]{geometry}
\usepackage{amsmath,amssymb}
\setlength{\parindent}{0pt}
\setlength{\parskip}{6pt}
\begin{document}
\section*{STAT 412 QZ08 --- Question 10}

\subsection*{Solution}

\[
H_0:\text{gender and whether a call is overturned are independent.}
\]
\[
H_1:\text{gender and whether a call is overturned are not independent.}
\]

Expected counts:
\[
E=\frac{(\text{row total})(\text{column total})}{\text{grand total}}.
\]
For example, the expected count for men with successful challenges is
\[
E=\frac{976(720)}{2402}=292.556.
\]
\[
\begin{array}{c|cc}
 & \text{Successful} & \text{Unsuccessful}\\ \hline
\text{Men} & 292.556 & 683.444\\
\text{Women} & 427.444 & 998.556
\end{array}
\]

The chi-square statistic is
\[
\chi^2=\sum\frac{(O-E)^2}{E}=10.988.
\]
For a \(2\times2\) table,
\[
df=(2-1)(2-1)=1.
\]

\[
\chi^2=10.988,\qquad df=1,\qquad p=0.0009.
\]

Since \(p<0.05\), reject \(H_0\). There is sufficient evidence to reject the claim of independence.

\subsection*{Final Answer}

\fbox{\begin{minipage}{0.94\linewidth}
\(\chi^2=10.988\), \(p=0.0009\). Reject \(H_0\); gender and challenge result are not independent.
\end{minipage}}

\end{document}
